\documentclass[12pt,a4paper]{article}
\usepackage[utf8]{inputenc}
\usepackage[T1]{fontenc}
\usepackage[french]{babel}
\usepackage{geometry}
\usepackage{titlesec}
\usepackage{parskip}
\usepackage{hyperref}
\usepackage{xcolor}
\usepackage{enumitem}
\usepackage{fancyhdr}
\usepackage{booktabs}
\usepackage{array}

\geometry{margin=2.5cm}

\definecolor{codeblue}{RGB}{0,70,150}
\definecolor{codered}{RGB}{180,0,0}
\definecolor{codegray}{RGB}{100,100,100}

\pagestyle{fancy}
\fancyhf{}
\rhead{Application de statistiques sportives}
\lhead{Fichier \texttt{\_\_main\_\_.py}}
\cfoot{\thepage}

\titleformat{\section}{\large\bfseries\color{codeblue}}{}{0em}{}[\titlerule]
\titleformat{\subsection}{\normalsize\bfseries\color{codered}}{}{0em}{}
\titleformat{\subsubsection}{\normalsize\itshape}{}{0em}{}

\newcommand{\fonction}[1]{\textbf{\texttt{#1}}}
\newcommand{\module}[1]{\textit{\texttt{#1}}}

\begin{document}

\begin{titlepage}
    \centering
    \vspace*{3cm}
    {\huge\bfseries Documentation technique\\[0.5em]}
    {\Large\bfseries Fichier \texttt{\_\_main\_\_.py}\\[2em]}
    {\large Application de statistiques sportives\\[1em]}
    {\normalsize Projet 1A — 2026\\[3em]}
    \hrule
    \vspace{1em}
    {\small Ce document décrit l'ensemble des fonctions et de la classe
    présentes dans le fichier \texttt{\_\_main\_\_.py}, qui constitue le
    point d'entrée de l'application et implémente l'intégralité de
    l'interface en ligne de commande.}
\end{titlepage}

\tableofcontents
\newpage

% ============================================================
\section{Architecture générale}

Le fichier \texttt{\_\_main\_\_.py} est organisé en trois parties
distinctes~:

\begin{enumerate}
    \item \textbf{Fonctions utilitaires de niveau module} — fonctions
    indépendantes pour l'affichage et la lecture du terminal.
    \item \textbf{Classe \texttt{CLI}} — contient toute la logique de
    navigation entre les pages de l'application.
    \item \textbf{Point d'entrée \texttt{main()}} — instancie la CLI
    et lance la boucle principale.
\end{enumerate}

La navigation repose sur un \textbf{pattern état-machine}~: chaque
\textit{page} est une fonction callable retournant soit une autre
fonction-page (avancer), soit la chaîne \texttt{"back"} (reculer),
\texttt{"quit"} (quitter) ou \texttt{None} (rester sur la page courante).
Deux piles \texttt{\_back} et \texttt{\_fwd} maintiennent l'historique,
à l'image d'un navigateur web.

% ============================================================
\section{Fonction \texttt{\_make\_sports\_config()}}

\textbf{Type~:} fonction de niveau module.

Construit et retourne le dictionnaire de configuration complet de
l'application. Tous les modules d'analyse (\texttt{src.Analysis.*})
sont importés \textbf{à l'intérieur} de cette fonction, de façon à
retarder leur chargement au premier appel (les données CSV ne sont
lues qu'à ce moment).

Chaque clé du dictionnaire est le nom affiché d'un sport. La valeur
associée décrit sa structure~:

\begin{itemize}
    \item \texttt{collectif} (bool) — détermine si le sport propose
    une vue Équipe et une vue Joueurs.
    \item \texttt{genre} (bool) — détermine si une sélection
    Hommes / Femmes est proposée en premier.
    \item \texttt{equipe} / \texttt{joueurs} / \texttt{stats} — listes
    de descripteurs de statistiques. Chaque descripteur contient~:
    \begin{itemize}
        \item \texttt{label} — texte affiché dans le menu.
        \item \texttt{fn} — la fonction d'analyse à appeler.
        \item \texttt{inputs} — liste de champs libres à saisir (vide
        dans la plupart des cas).
        \item \texttt{selector} (optionnel) — décrit comment générer
        la liste de choix. Deux variantes existent~:
        \begin{itemize}
            \item \texttt{options\_fn} — liste simple, affichée en
            bloc (ex~: liste d'équipes).
            \item \texttt{options\_df\_fn} — DataFrame avec pagination
            et filtres (ex~: liste de joueurs).
        \end{itemize}
    \end{itemize}
\end{itemize}

% ============================================================
\section{Fonctions utilitaires d'affichage}

Ces six fonctions sont de simples helpers de terminal, sans état.

\subsection{\texttt{\_titre(texte)}}
Affiche un titre encadré par une ligne de tirets (\texttt{────}).
Utilisé en en-tête de chaque page.

\subsection{\texttt{\_option(num, texte)}}
Affiche une ligne de menu au format \texttt{[N]  texte}.

\subsection{\texttt{\_nav(peut\_reculer, peut\_avancer)}}
Affiche la barre de navigation contextuelle en bas de page~:
\texttt{[-] Retour}, \texttt{[+] Suivant}, \texttt{[q] Quitter}.
Les options retour et suivant sont conditionnées par les paramètres
booléens.

\subsection{\texttt{\_lire(prompt)}}
Lit une saisie utilisateur via \texttt{input()} et retourne la chaîne
nettoyée (\texttt{strip()}). Le prompt par défaut est
\texttt{"Votre choix~:"}.

\subsection{\texttt{\_afficher\_resultat(result)}}
Affiche un résultat~: si \texttt{result} est un \texttt{DataFrame}
non vide, il est converti en texte via \texttt{to\_string(index=False)}.
Si le DataFrame est vide, affiche \texttt{"(aucun résultat)"}. Si la
fonction appelée a déjà imprimé directement (ex~: bracket), cette
fonction ne fait rien.

\subsection{\texttt{\_extract\_year(val)}}
Extrait une année depuis des formats hétérogènes~:
\begin{itemize}
    \item \texttt{datetime} ou \texttt{Timestamp} → attribut \texttt{.year}.
    \item \texttt{int} ou \texttt{float} supérieur à 19\,000\,000 →
    format \texttt{YYYYMMDD}, division par 10\,000.
    \item \texttt{int} ou \texttt{float} inférieur → déjà une année.
    \item \texttt{str} → lecture des 4 premiers caractères.
\end{itemize}
Cette fonction est utilisée par le système de filtres sur les années de
naissance, où les données proviennent de sources aux formats variés.

% ============================================================
\section{Classe \texttt{CLI}}

La classe \texttt{CLI} encapsule l'intégralité de la navigation de
l'application. Elle maintient~:
\begin{itemize}
    \item \texttt{\_back} — pile des pages visitées (navigation arrière).
    \item \texttt{\_fwd} — pile des pages en avant (navigation avant).
    \item \texttt{\_filters} — dictionnaire des filtres actifs
    (\texttt{type\_sport}, \texttt{categorie}, \texttt{type\_competition}).
    \item \texttt{\_config} — la configuration complète des sports,
    chargée au démarrage.
\end{itemize}

\subsection{Méthode \texttt{\_\_init\_\_()}}
Initialise les deux piles de navigation à des listes vides et les
filtres à \texttt{None}. La configuration \texttt{\_config} est
initialisée à \texttt{None} et sera remplie lors de l'appel à
\texttt{run()}.

\subsection{Méthode \texttt{run()}}
Point de départ de la boucle principale. Charge la configuration via
\texttt{\_make\_sports\_config()}, puis entre dans une boucle infinie
qui~:
\begin{enumerate}
    \item Appelle la page courante (\texttt{current()}).
    \item Interprète le résultat~:
    \begin{itemize}
        \item \texttt{"quit"} → affiche le message d'au revoir et sort.
        \item \texttt{"back"} → dépile \texttt{\_back}, empile dans
        \texttt{\_fwd}, change de page.
        \item \texttt{"forward"} → dépile \texttt{\_fwd}, empile dans
        \texttt{\_back}, change de page.
        \item Un \textit{callable} → empile la page courante dans
        \texttt{\_back}, vide \texttt{\_fwd}, change de page.
        \item \texttt{None} → reste sur la page courante (réaffichage).
    \end{itemize}
\end{enumerate}

\subsection{Méthode \texttt{\_sports\_visibles()}}
Appelle \texttt{src.Model.sport.filtrer()} avec les trois filtres actifs
pour obtenir la liste des sports correspondants, puis la croise avec
les clés de \texttt{\_config} pour ne retourner que les sports
effectivement disponibles dans l'application.

\subsection{Méthode \texttt{\_filtres\_str()}}
Retourne une chaîne lisible des filtres actifs pour l'affichage dans
le menu principal (ex~: \texttt{"Collectif, E-sport"}). Retourne
\texttt{"Aucun"} si aucun filtre n'est activé.

% ============================================================
\section{Pages de navigation de la classe \texttt{CLI}}

\subsection{Méthode \texttt{\_page\_accueil()}}
Page racine de l'application. Affiche les filtres actifs et la liste
numérotée des sports visibles. L'option \texttt{[0]} mène à la page
des filtres, les autres options mènent aux pages des sports. Retourne
une fonction-page ou une chaîne de contrôle.

\subsection{Méthode \texttt{\_page\_filtres()}}
Page de configuration des filtres. Propose trois catégories de filtres
combinables~:
\begin{itemize}
    \item \textbf{Format}~: Collectif / Individuel.
    \item \textbf{Discipline}~: Sport / E-sport.
    \item \textbf{Compétition}~: Par points / Éliminatoire / Mixte.
\end{itemize}
Chaque catégorie dispose d'une option de réinitialisation. Les filtres
actifs sont indiqués par le symbole \texttt{[✓]} en face de l'option.
La modification d'un filtre retourne \texttt{None} (la page se
réaffiche immédiatement avec l'état mis à jour).

% ============================================================
\section{Factories de pages de la classe \texttt{CLI}}

Ces méthodes ne sont pas elles-mêmes des pages~: elles \textbf{créent
et retournent} des fonctions-pages (pattern \textit{factory}). Ce
choix permet de paramétrer dynamiquement le comportement de chaque
page sans dupliquer de code.

\subsection{Méthode \texttt{\_make\_page\_sport(sport)}}
Lit la configuration du sport (\texttt{collectif}, \texttt{genre}) et
délègue vers la factory appropriée~:
\begin{itemize}
    \item Avec genre → \texttt{\_make\_page\_genre}.
    \item Sans genre, collectif → \texttt{\_make\_page\_cat}.
    \item Individuel → \texttt{\_make\_page\_stats}.
\end{itemize}

\subsection{Méthode \texttt{\_make\_page\_genre(titre, config)}}
Crée une page proposant le choix Hommes / Femmes. Selon la valeur de
\texttt{collectif}, délègue ensuite vers \texttt{\_make\_page\_cat}
(sport collectif avec genre, ex~: Volleyball) ou
\texttt{\_make\_page\_stats} (sport individuel avec genre, ex~: Tennis).

\subsection{Méthode \texttt{\_make\_page\_cat(titre, stats\_equipe, stats\_joueurs)}}
Crée une page proposant le choix Équipe / Joueurs. Chaque choix mène
vers une page \texttt{\_make\_page\_stats} avec la liste de stats
correspondante.

\subsection{Méthode \texttt{\_make\_page\_stats(titre, stats)}}
Crée une page listant toutes les statistiques disponibles pour un sport
ou une catégorie. Chaque option est un descripteur de stat issu de
\texttt{\_make\_sports\_config()}.

\subsection{Méthode \texttt{\_make\_page\_stat(stat)}}
Crée la page d'exécution d'une statistique. Si le descripteur contient
un sélecteur de type \texttt{options\_df\_fn} (DataFrame paginé),
délègue vers \texttt{\_make\_page\_selector\_paged}. Sinon, gère~:
\begin{enumerate}
    \item Le sélecteur simple (\texttt{options\_fn}) — affiche une liste
    numérotée et attend un choix.
    \item Les champs libres (\texttt{inputs}) — demande une saisie
    pour chaque paramètre.
    \item L'exécution de la fonction d'analyse (\texttt{fn(**kwargs)}).
    \item L'affichage du résultat via \texttt{\_afficher\_resultat}.
    \item Le transfert vers \texttt{\_attendre\_apres\_resultat}.
\end{enumerate}

\subsection{Méthode \texttt{\_make\_page\_selector\_paged(sel, stat)}}
Crée une page de sélection paginée (20 éléments par page) avec filtres.
C'est la page la plus complexe de l'application. Elle maintient un
dictionnaire d'état local \texttt{state} persistant entre les appels~:
\begin{itemize}
    \item \texttt{page} — index de la page courante.
    \item \texttt{filters} — filtres actifs sur le DataFrame.
    \item \texttt{df} — DataFrame chargé une seule fois.
    \item \texttt{\_filter\_applied} — flag de retour automatique après
    application d'un filtre.
\end{itemize}

La sous-fonction \texttt{\_apply\_filters(df)} applique les filtres
de trois types~:
\begin{itemize}
    \item \textbf{search} (\texttt{\_\_search\_\_}) — recherche par
    préfixe sur chaque mot du nom affiché.
    \item \textbf{category} — égalité exacte sur une colonne.
    \item \textbf{year} (\texttt{\_\_year\_\_<col>}) — filtre par année
    exacte ou plage d'années (ex~: \texttt{1990-2000}), via
    \texttt{\_extract\_year}.
\end{itemize}

Les commandes disponibles sont \texttt{[<]}/\texttt{[>]} pour la
pagination, \texttt{[f]} pour ouvrir le menu de filtres, \texttt{[r]}
pour réinitialiser les filtres, et un numéro pour sélectionner un
élément.

\subsection{Méthode \texttt{\_make\_page\_filter(state, filters\_cfg, df)}}
Crée la page listant les filtres disponibles pour un sélecteur paginé.
Gère les trois types de filtres~:
\begin{itemize}
    \item \textbf{search} — saisie libre de lettres initiales.
    \item \textbf{category} — délègue vers
    \texttt{\_make\_page\_filter\_category}.
    \item \textbf{year} — saisie d'une année ou d'une plage
    (\texttt{1990} ou \texttt{1990-2000}). Valide le format et lève
    un message d'erreur si invalide.
\end{itemize}
Après application d'un filtre, le flag \texttt{\_filter\_applied} est
levé, ce qui déclenche un retour automatique vers le sélecteur au
prochain cycle.

\subsection{Méthode \texttt{\_make\_page\_filter\_category(state, fconf, df)}}
Crée une page paginée (20 valeurs par page) pour sélectionner une
valeur de filtre catégoriel. Calcule les valeurs uniques de la colonne
filtrée, les trie et les pagine. La sélection d'une valeur met à jour
\texttt{state["filters"]}, lève le flag \texttt{\_filter\_applied}
et retourne \texttt{"back"} pour remonter automatiquement vers le
sélecteur.

% ============================================================
\section{Méthodes de fin de parcours de la classe \texttt{CLI}}

\subsection{Méthode \texttt{\_attendre\_apres\_resultat(result, stat\_label)}}
Boucle d'attente affichée après l'exécution d'une statistique.
Propose~:
\begin{itemize}
    \item \texttt{[e] Exporter CSV} — uniquement si \texttt{result}
    est un \texttt{DataFrame} non vide.
    \item \texttt{[+] Suivant} — uniquement si la pile \texttt{\_fwd}
    n'est pas vide.
    \item \texttt{[-] Retour} et \texttt{[q] Quitter} — toujours
    présents.
\end{itemize}
La méthode boucle indéfiniment jusqu'à ce que l'utilisateur choisisse
une action de navigation.

\subsection{Méthode \texttt{\_exporter\_csv(df, stat\_label)}}
Sauvegarde un \texttt{DataFrame} dans un fichier CSV. Génère un nom
de fichier par défaut en combinant un slug du label de la stat
(caractères alphanumériques uniquement via regex
\texttt{[{\textasciicircum}a-z0-9]+}) et un horodatage au format
\texttt{YYYYMMDD\_HHMMSS}. L'utilisateur peut accepter ce nom ou en
saisir un autre. L'encodage utilisé est \texttt{utf-8-sig} (compatible
avec Excel).

% ============================================================
\section{Fonction \texttt{main()}}

Point d'entrée de l'application. Instancie la classe \texttt{CLI} et
appelle \texttt{run()}. Cette fonction est appelée depuis le bloc
\texttt{if \_\_name\_\_ == "\_\_main\_\_"}, ce qui permet aussi
d'importer le module sans lancer l'interface.

% ============================================================
\section{Synthèse du flux de navigation}

\begin{center}
\begin{tabular}{>{\ttfamily}p{4.5cm} p{9cm}}
\toprule
\textbf{Page / Factory} & \textbf{Rôle} \\
\midrule
\_page\_accueil & Menu principal avec liste des sports filtrés \\
\_page\_filtres & Configuration des filtres Sport/Format/Compétition \\
\_make\_page\_sport & Aiguillage selon le type de sport \\
\_make\_page\_genre & Sélection Hommes / Femmes \\
\_make\_page\_cat & Sélection Équipe / Joueurs \\
\_make\_page\_stats & Liste des statistiques disponibles \\
\_make\_page\_stat & Exécution d'une statistique (sélecteur simple) \\
\_make\_page\_selector\_paged & Sélection paginée avec filtres (joueurs) \\
\_make\_page\_filter & Menu des filtres disponibles \\
\_make\_page\_filter\_category & Sélection paginée d'une valeur catégorielle \\
\_attendre\_apres\_resultat & Navigation et export après affichage d'un résultat \\
\bottomrule
\end{tabular}
\end{center}

\begin{center}
\begin{tabular}{>{\ttfamily}p{4cm} p{9.5cm}}
\toprule
\textbf{Valeur retournée} & \textbf{Effet dans la boucle \texttt{run()}} \\
\midrule
callable & Avancer vers la nouvelle page, empiler la courante dans \_back \\
"back" & Reculer d'une page (dépiler \_back, empiler dans \_fwd) \\
"forward" & Avancer d'une page (dépiler \_fwd, empiler dans \_back) \\
"quit" & Terminer l'application \\
None & Rester sur la page courante (réaffichage) \\
\bottomrule
\end{tabular}
\end{center}

\end{document}
